%% kapa: exam=1 round=36 year=2025 period=1 form=A subject=economics subject_ko=경제학원론 seq=5 number=45
%% source: Q-Net 감정평가사 기출문제 (cst00309) artlSeq=5231525 file=raw/1008/5231525/2025년 제36회 감정평가사 1차시험 1교시 문제지 원본.pdf
%% conversion: kordoc 4.24.51 -> tools/gen1.py (auto); review=proofread-sonnet; flags=none
\begin{question}{45}
시장실패에 관한 설명으로 옳지 않은 것은?
\begin{choices}
  \item 공유자원의 비극은 소유권이 불명확하여 자원이 과대하게 사용되는 문제이다.
  \item 부정적 외부성은 사회적 후생손실을 발생시키지만, 긍정적 외부성은 사회적 후생손실을 발생시키지 않는다.
  \item 정보의 비대칭성은 역선택과 도덕적 해이를 유발할 수 있다.
  \item 공공재는 비배제성의 특징으로 인해 무임승차자 문제가 발생한다.
  \item 시장실패에 대응하기 위한 정부개입이 오히려 효율성을 저하시킬 수 있다.
\end{choices}
\end{question}
\answer{2}
